%!TEX root = ../main.tex
%% English abstract — task → limits of prior definitions → contribution → benefits → experiment summary. Skirmish external sentence: author decision A2.
\begin{abstracten}
Before a fight starts in League of Legends, how much of its consequence is visible in public match state? This thesis studies the problem of predicting, for engagements constructed from kill records, the direction of change in estimated match win probability across the engagement from the public state available at an operational cutoff before the first kill. Public telemetry records state once a minute and events to the millisecond without damage events, so engagement boundaries must be reconstructed from kills; the effect of a fight on the match outcome is never observed; and the current win probability and the clock already explain much of what follows. Earlier engagement-level prediction set its constants by game reasoning and targeted a fight-local exchange-value label.

The thesis connects, as one procedure, the construction of engagement instances from asynchronously recorded public data, the definition of a match-linked outcome value, prediction from pre-fight information, and the verification of what that prediction means and where it holds. Engagements are anchored on kills, but the temporal boundary (\constG{} s) and the spatial boundary (\constD{} units) are estimated from the corpus, the presence-gate radius (\constR{} units) and lead (\constB{} s) are fixed at values the game itself applies, and every constant that shapes the engagement set is tabulated with its source class. On \numMatchesTotal{} Korean-server matches from three patches, detected engagements are grouped by the smaller side's participation into teamfights (at least four champions) and skirmishes (two or three). The outcome value is the direction of change in estimated match win probability across the engagement interval under an evaluator trained on match outcomes and then frozen (strategic value improvement, SVI); the evaluator's quality, the separation of direction, mean and scale, a contrast with matched quiet intervals, correspondence with material outcomes and with alternative outcome definitions, label stability across horizons, evaluators and stopping rules, and an arithmetic decomposition into clock and observation-update components bound what this label measures.

Within each cohort, a regularized logistic model on pre-fight state lowers the match-weighted Brier score relative to a spline baseline in initial win probability and game time by \dBrierT{} (teamfights) and \dBrierS{} (skirmishes) on a held-out patch, with match-cluster bootstrap intervals excluding zero. The teamfight model carries part of its signal to skirmishes, and training on both cohorts together did not improve either cohort in the setting examined. In an information-group comparison under one learning procedure, adding event features to frame features improved the score in both cohorts, while the gain from adding frame features to event features differed between cohorts. Labels were largely preserved across neighbouring evaluators and stopping rules, but the sign of the prediction contrast reversed for small changes. The lift is smaller in balanced states, survives filters on small changes, coincides with a larger discrimination component, and keeps its sign in cases with a fresh observation frame. In the external samples evaluated (a later patch, two regions) the teamfight advantage was not observed, and the skirmish advantage was uncertain at the family-corrected level. The two cohorts are reported separately. The results concern a model-defined SVI conditional on the frozen evaluator and the engagement definition, and they state the additional predictive value that public state provides together with the information composition, evaluator, stopping rule and external environment on which that value depends.
\end{abstracten}
